\section{Introduction}\label{sec:intro}
Information disclosure is part of the operation of modern power infrastructure. Operators publish forecasts, prices, reserve requirements and schedules so that market participants and other stakeholders can coordinate their decisions. At the same time, unit-specific information and some actual operating quantities are restricted or disclosed only after a delay. Protecting these quantities while releasing useful information is an ongoing infrastructure-management task.

The release catalogue is organized into fields, but the underlying system is interconnected. Loads, generation, network constraints and prices describe different aspects of the same operating state. Information that is weakly revealing in isolation may therefore become revealing when combined with another published series. A price at a generating unit's bus illustrates the issue: without the unit's output, it may say little about a private offer parameter; with the output, it can identify that parameter through the active offer segment. Congestion can change the relation for prices measured elsewhere in the grid. The risk of the release consequently depends on both the field and the background information available to its recipient.

This problem complements the confidentiality, integrity and availability concerns addressed by infrastructure protection. Recent power-system studies have examined how corrupted smart-meter measurements affect state estimation and protective measures~\cite{farooq2024grid}, and how data-driven recovery can preserve load-frequency control under disrupted communications~\cite{syrmakesis2024darlfc}. Here the information is released through the intended disclosure process. The protection task is to determine what sensitive quantities can be inferred from those releases and which combinations create the exposure.

Disclosure rules already provide a practical starting point. They distinguish field categories and release times, for example forecasts before clearing, prices after clearing and actual quantities on the following day~\cite{nea2024disclosure}. Classification standards require the effects of aggregation and processing to be considered~\cite{gbt43697}. However, an individual field's predictive relationship with a target does not describe what it adds in the presence of other fields. An operational audit needs a field-level result accompanied by evidence of these dependencies.

We propose such a result: a profile consisting of single-field risk, background amplification and high-gain background combinations. The profile is derived from the out-of-sample inferability of a target using a specified set of fields. Its principal measure is the largest marginal gain from adding the field to any background within a chosen budget, relative to information already public. This budgeted marginal risk gives a consistent field-level measurement, while the associated backgrounds identify the release combinations that require attention.

The audit must also be repeatable as information is released or reclassified. We share a learned representation across the many field sets and fit a small readout for each set, reducing the cost of scanning to tens of milliseconds per query. Dedicated retraining then verifies selected backgrounds. Fig.~\ref{fig:framework} shows how these steps support a release review.

The contributions are threefold. First, we connect field-level grading to measurable set-level inferability and establish its relation to threshold crossings and minimal unsafe combinations. Second, we explain a controlled disclosure mechanism through power dispatch, unit offers and congestion, and evaluate the prevalence of amplification across four power-system data sets. Third, we demonstrate how verified field profiles and combination structure can support withholding and repeated review at practical computational cost. The central outcome is an interpretable account of when a field becomes informative, suitable for incorporation into an operator's existing disclosure process.

\begin{figure}[!tb]
\centering\includegraphics[width=\textwidth]{fig1_framework.pdf}
\caption{Background-aware review of power-system information. Coupled operational fields create context-dependent inferability. The audit reports the single-field contribution, its amplification and supporting combinations, using an efficient scan followed by retraining of selected backgrounds. The network and miniature profile are illustrative.}\label{fig:framework}
\end{figure}
